\subsection{Restoring the complete odd root}
\label{sec:odd-complete:assembly}

The terminal directions in Table~\ref{tab:odd186:directions}
are shared by the two root classes, apart from the short-one
direction \eqref{eq:odd-short:directions}. The preceding
complete-certificate proofs give the following global rates:
\begin{equation}\label{eq:odd-complete:catalogue}
\begin{array}{c|c|c|l}
 \text{terminal}&pW_a\bar f_b&f_a h_bf_c&\text{proof}\\\hline
 (1,t)&705/128&11/2&\text{whole-message short-one packets}\\
 (2,t),\ t\ge4&6&23/4&\text{two-arm theorem and short exceptions}\\
 (2,3)&6&6&\text{two-arm theorem and short exceptions}\\
 (3,t),\ t\ge4&6&6&\text{Theorem~\ref{thm:odd-three:roots}}\\
 \text{even lists of minimum at least four}&8&11/2
          &\text{Proposition~\ref{prop:odd-even:roots}}\\
 \text{other even multipair lists}&11/2&11/2
          &\text{Proposition~\ref{prop:odd-even:roots}}\\
 \text{odd lists of minimum one}&705/128&11/2
          &\text{Section~\ref{sec:odd-complete:q-terminals}}\\
 \text{odd lists of nonunit minimum}&13/2&6
          &\text{Section~\ref{sec:odd-complete:inward-anchor}}.
\end{array}
\end{equation}
The positive-root entry in the first row means its actual special
direction \eqref{eq:odd-short:directions}, including all later
negative pairs. Every other positive root uses $f_a h_bf_c$.

\begin{lemma}[The full low-domain vector]
\label{lem:odd-complete:vector-cap}
On $0<d\le1/64$, each base certificate used in
\eqref{eq:odd-complete:catalogue} has its original pointwise form
\[
 \log H_0\le a_0+BV_0+W_0,
 \qquad\|V_0\|<(7/2)\sqrt d,
\]
with its already paid complete fourth exponent at most
$-27d/5$ at a minimum-one root and $-11d/2$ at a positive root.
Here $W_0=0$ or the same paid unit-arm curvature term.
\end{lemma}
\begin{proof}
The minimum-one short-one bound is
\eqref{eq:odd186:short-one-certificate}, with vector norm
$3\sqrt d/2$. Its two-arm-minimum-two counterpart is supplied
by Theorem~\ref{thm:odd-two:roots}, with norm at most $2\sqrt d$.
For the other minimum-one rows the transferred nonconstant root
precursor has norm at most $\sqrt d/2$; the original constant
or curved coefficient is the one used in its proof.

For the positive root outside its short-one special direction the
complete vector is exactly
\[
 V_0=Q(B)u+G_{a-1}(-B)u+V_{bc}+\gamma\one.
\]
The two root precursors have norms below
$(6/5)\sqrt d$ and $(3/8)\sqrt d$.
The following bounds for the core
$(\|Qu\|+|\gamma|)/\sqrt d$ include the actual constant:
\[
\begin{array}{c|c}
 q\text{-first-row and retained-square packets}&5/3\\
 \text{bare-exponent packets, first three windows}
          &1,\ 9/5,\ 147/80\\
 \text{bare-exponent packets, long window}&167/120\\
 \text{general inward anchor}&29/20\\
 \text{short inward anchor, first three windows}
          &\begin{gathered}9/5,\quad1+(13/20)(3/5)+1/16,\\
             1+1/128+(13/32)(3/5)+1/4\end{gathered}\\
 \text{short inward anchor, long window}&7/5.
\end{array}
\]
For the exponent rows use $\rho\le1$,
$|\delta|\le6d/5$, $|\epsilon|\le17d/10$ and $h\le1$;
at $L=3$ the genuine $G_2$ contribution is $6/5$, the scalar
term in $Q$ costs at most $3/160$, the attenuated remainder
at most $13/32$, and the constant vector at most $17/80$.
The general anchor keeps $\mu_A\le1$ and its $G_2$ exception.
The short anchor uses its two actual $D$ values and
\eqref{eq:odd-anchor:short-rows}--\eqref{eq:odd-anchor:multiple-long};
its $|D-d|u$ term is distinct from $\gamma\one$.
All entries are less than $19/10$, so
$\|V_0\|<(19/10+6/5+3/8)\sqrt d=(139/40)\sqrt d$.
The special short-one direction was paid directly for all its
root pairs in \eqref{eq:odd-short:whole-pairs}.
\end{proof}

Let $k\ge1$ further ordered unequal root pairs follow the base
root triple. At a minimum-one root they are $W_s\bar f_t$;
at a positive root they are $h_sf_t$.
Their disjoint original exponent intervals give one finite mask,
with scalar loss $16kd/5$, vector norm below $3\sqrt d/8$ and
signed mean at most $d/12$. In the positive direction its positive
mean powers first occur at original power eight, giving
$\delta_4/(1-\theta^2)<d/12$; complete reversal gives the
minimum-one direction. Applying
Lemma~\ref{lem:odd-common:extra-pairs} to the full vector above
gives the additional rate
\begin{equation}\label{eq:odd-complete:later-rate}
 \zeta_k=2737/3840+64(k-1)/5.
\end{equation}
The high-domain estimate in that lemma uses these same products.
For the short-one positive root, the directly proved increment
$64k/5$ in \eqref{eq:odd-short:whole-pairs} is stronger.

Genuine equal pairs are oriented as $f_\ell h_\ell$ or their
complete reverse. Their products are at most one. Positivity
therefore restores all root and terminal decorations in the
actual message. Finally apply the unanchored ordinary outlet
\eqref{eq:odd-common:outlet}, independently of the inner $L$.
The base losses $27/5$ and $11/2$ give the respective output
losses $6/95$ and $5/19$ for all $M\ge1$, and each later root
pair supplies twice its additional fourth loss.
